\documentclass[journal]{IEEEtran}

\usepackage{amsmath}
\usepackage{booktabs}
\usepackage{graphicx}
\usepackage{siunitx}

\title{Sampler-Aperture-Aware Implementation Analysis of FPGA Ring-Oscillator TRNGs}

\author{TODO}

\begin{document}

\maketitle

\begin{abstract}
TODO. This TVLSI-track manuscript is intentionally separate from the TIM-track manuscript. Current claims must be filled only after the offline model and follow-on validation evidence mature.
\end{abstract}

\begin{IEEEkeywords}
Ring oscillator, true random number generator, FPGA, physical implementation, timing, sampling aperture.
\end{IEEEkeywords}

\section{Introduction}
TODO. Frame the problem as an implementation-methodology problem rather than only a measurement anomaly.

\section{Background and Motivation}
TODO. Define restarted RO-TRNG sampling, XOR aggregation, warmup, placement/routing sensitivity, and the sampler-aperture hypothesis.

\section{Stochastic Contributor and Sampler-Aperture Model}
TODO. Introduce contributor variables, signed bias factors, XOR cancellation, board and implementation variables, and route/timing feature variables.

\section{Implementation and Evidence Extraction Methodology}
TODO. Describe reduced-XOR decomposition, repeat protocol, route/PIP/net-delay audit, and board-level held-out sampler checks.

\section{Results}
TODO. Insert generated TVLSI offline model tables only after claim-evidence review.

\section{Discussion and Design Implications}
TODO. Separate observed effects from candidate design rules.

\section{Limitations and Future Work}
TODO. Calibrated jitter propagation, metastability transfer fitting, formal timing-to-aperture derivation, and broader PVT/toolflow validation belong here unless directly validated.

\section{Conclusion}
TODO.

\bibliographystyle{IEEEtran}
\bibliography{../refs/references}

\end{document}
